\documentclass[11pt]{article}
\usepackage[margin=2.6cm]{geometry}
\usepackage[T1]{fontenc}
\usepackage[utf8]{inputenc}
\usepackage{lmodern}
\usepackage{amsmath,amssymb}
\usepackage{array,longtable}
\usepackage{listings}
\usepackage{xcolor}
\usepackage[hidelinks]{hyperref}
\usepackage{parskip}

\lstset{
  basicstyle=\ttfamily\small,
  breaklines=true,
  columns=fullflexible,
  keepspaces=true,
  frame=single,
  framerule=0.2pt,
  xleftmargin=0.6em,
  aboveskip=0.7em,
  belowskip=0.7em,
  literate={—}{{---}}1 {°}{{\ensuremath{^\circ}}}1 {≈}{{\ensuremath{\approx}}}1
           {←}{{\ensuremath{\leftarrow}}}1 {→}{{\ensuremath{\rightarrow}}}1
           {├}{{|--}}3 {└}{{`--}}3 {│}{{|}}1 {─}{{-}}1 {▶}{{\ensuremath{\triangleright}}}1
           {⏸}{{||}}2 {■}{{\ensuremath{\blacksquare}}}1 {◀}{{\ensuremath{\triangleleft}}}1
           {⇠}{{\ensuremath{\dashleftarrow}}}1 {⇢}{{\ensuremath{\dashrightarrow}}}1
           {⌂}{{home}}4 {×}{{\ensuremath{\times}}}1 {↑}{{\ensuremath{\uparrow}}}1
           {↓}{{\ensuremath{\downarrow}}}1 {−}{{-}}1 {↔}{{\ensuremath{\leftrightarrow}}}1
}

\newcommand{\kw}[1]{\texttt{\uppercase{#1}}}
\newcommand{\code}[1]{\texttt{#1}}

\title{The \code{posim} Graphical Scene Window\\[2pt]
       \large Research, Design \& Complete Guide (\code{scene\_info})}
\author{physical\_object\_simulator}
\date{July 22, 2026}

\begin{document}
\maketitle

\noindent\emph{This is the typeset twin of \code{scene\_info.md} ---
companion to \code{grammar.pdf} and
\code{physical\_object\_simulator.pdf}.}

This document explains, for a reader with \textbf{zero prior knowledge}:
\begin{enumerate}
\item what the graphical scene window is and how to drive it,
\item the \textbf{research} performed on seven established physics
      simulators before designing it,
\item a \textbf{comparison chart and concept tree} of that ``scene
      science'',
\item which ideas we recommend (and chose) to \textbf{port to pure
      Rust}, and
\item the full technical documentation of the implementation, its
      asynchronous notebook link, and how every claim was
      \textbf{verified}.
\end{enumerate}

\tableofcontents

\section{What is the scene window?}

\code{posim} is a physics simulator you drive by typing commands
(\code{NEW SPHERE \{ ... \}}, \code{STEP 0.1}, \ldots). Numbers are
exact but hard to \emph{picture}. The \textbf{scene window} is a live
3-D view of every object in your simulation that opens \textbf{in your
web browser, outside the notebook}, the moment you type:

\begin{lstlisting}
scene create
\end{lstlisting}

You get back something like:

\begin{lstlisting}
scene window created: http://127.0.0.1:41234/
(asked your desktop to open it; if no window appeared, open that address yourself)
showing 4 entities; SCENE START begins the evolution - HELP lists all scene commands
\end{lstlisting}

The middle line reports what was actually attempted rather than assuming
it worked. posim tries exactly one command, \kw{xdg-open} \textit{url},
which is a Linux/BSD utility. When that launch does not happen --- on
macOS and Windows, where \kw{xdg-open} does not exist, or whenever
\verb|$POSIM_NO_BROWSER| is set --- the line reads instead:

\begin{lstlisting}
(no browser was launched - open that address yourself)
\end{lstlisting}

In both cases the URL is printed and the window works; paste it into any
browser (\kw{open} \textit{url} on macOS, \kw{start} \textit{url} on
Windows). The listener is bound to \texttt{127.0.0.1}, so it is not
reachable from another machine.

If \textbf{no} bodies exist yet, the reply says so instead --- the
window would be an empty grid --- and reminds you that the scene draws
rigid bodies only (\kw{new} adds them, then \kw{scene} \kw{refresh});
quantum problems are viewed with \kw{qm} \kw{animate} /
\kw{qm2} \kw{animate}, not the scene window.

The window shows \textbf{all simulator entities} (points, spheres,
cuboids, tori, disks, cylinders, dumbbells), a ground grid, the world
axes, motion trails, and object labels (an object registered with
\kw{new} \code{...} \kw{as} \code{<name>} is labeled by that name).
If the notebook created a rigid
bounding box (\kw{box} \code{<size>}), the window draws its interior
as a dashed light-blue wireframe --- the six infinitely-massive wall
slabs that implement it are represented by that wireframe, never
drawn as six giant cuboids. The window keeps itself up to date over a
private network link to the simulator --- you never copy data by
hand. It can also \emph{talk back}: errors and user actions inside
the window flow to the notebook asynchronously (\S\ref{sec:async}).

No extra software is needed: the whole server is built into
\code{posim} itself using only the Rust standard library --- zero
external dependencies, zero \code{unsafe} code.

\subsection{The toolbar (top of the window)}

\begin{center}
\begin{tabular}{lp{9.2cm}}
\textbf{control} & \textbf{what it does}\\\hline
Start & begin time-stepped evolution, forward in time\\
Pause & freeze the evolution (resume with Start)\\
Stop & halt evolution \emph{and clear the recorded history}\\
Reverse & play \textbf{backward} through the recorded history\\
Reset & re-initialize the playback: \textbf{every value and the time return to their initial values}; Start then runs the simulation again from the beginning\\
step back / step & single step backward / forward\\
dt \code{\_\_\_\_} Set & change the playback time step\\
zoom + / zoom -- & zoom in / out\\
View & reset the camera to its home position\\
Grid / Trails / Labels & toggle the ground grid, motion trails, object labels\\
Contacts & toggle contact-normal arrows drawn at collision points\\
? & show the controls cheat-sheet\\
\end{tabular}
\end{center}

\subsection{The ``conserved quantities'' readout (top-left)}

A permanent labeled overlay (just below the toolbar) shows the three
conserved quantities of the playback copy, updated live at every
frame: \textbf{E} (total energy), \textbf{P} (total linear momentum)
and \textbf{L} (total angular momentum about the origin). P and L are
printed as their three components plus the magnitude \code{|.|}.
Watching E, P and L sit still through a collision \emph{is} the
physics lesson --- the values come straight from the frame protocol's
\code{energy}, \code{p} and \code{l} fields (\S5.1), computed on the
playback copy every tick.

\subsection{The status bar (bottom of the window)}

A live one-line dashboard: \textbf{connection dot} (green = connected
to posim), playback \textbf{mode} (\code{stopped} / \code{running} /
\code{paused} / \code{reversing}), simulated time \textbf{t}, time
step \textbf{dt}, total energy \textbf{E}, \textbf{bodies} count,
\textbf{hidden} count, \textbf{contacts} count (of the latest playback
step), \textbf{history} (frames available for
Reverse), the \textbf{camera} yaw/pitch/distance, and the rendering
\textbf{fps}.

\subsection{Controls (mouse \& keyboard) --- all verified in a real
browser}

\begin{center}
\begin{tabular}{lp{8.6cm}}
\textbf{gesture} & \textbf{effect}\\\hline
left / right arrow keys & \textbf{translate} the view left / right\\
up / down arrow keys & translate the view up / down\\
left-click + drag & \textbf{rotate} (orbit) around the scene\\
mouse wheel & \textbf{zoom} in / out\\
\code{+} / \code{-} keys & zoom in / out (documented keyboard shortcut)\\
Shift+drag or right-drag & translate (pan) with the mouse\\
Space & start / pause playback\\
\code{R} & reset the view\\
\code{G} / \code{T} / \code{L} & toggle grid / trails / labels\\
\code{C} & toggle contact-normal arrows\\
\code{H} or \code{?} & toggle the help overlay\\
\end{tabular}
\end{center}

\section{The research: seven simulators investigated}

Before writing a line of code we sent research agents to study how
seven established, widely-used simulators present and control a
graphical scene. A summary of each, then the comparison chart.

\subsection{PyChrono / Project Chrono (C++, Python bindings)}

Multi-physics engine used in engineering and robotics. Rendering is
\textbf{in-process} through an abstract \emph{visual system} layer
with pluggable backends --- Irrlicht
(\code{ChVisualSystemIrrlicht}), VulkanSceneGraph
(\code{ChVisualSystemVSG}), raw OpenGL, plus offline Blender/POVRay
export. Typical loop: \code{AttachSystem(sys)}, \code{Initialize()},
then \code{Run()/BeginScene()/Render()/EndScene()} around
\code{ChSystem::DoStepDynamics(dt)} --- \emph{you} own the loop and
the time step; pause/step are gates you place around
\code{DoStepDynamics}. Scene decoration: \code{AddCamera(pos,
target)}, \code{AddTypicalLights()}, \code{AddSkyBox()},
\code{EnableShadows()}. Viewer: left-drag rotate, right-drag pan,
wheel zoom, \code{i} info panel, Space pause.

\subsection{PyBullet (C++ engine, Python module)}

The gold standard in robotics labs. \textbf{Client--server by
design}: a client sends \emph{command structs}, the physics server
returns \emph{status structs} --- identical code whether the transport
is same-process (\code{connect(DIRECT)}/\code{connect(GUI)}), shared
memory, UDP, or TCP. Key calls: \code{stepSimulation()},
\code{setTimeStep(dt)}, \code{setRealTimeSimulation(0/1)},
\code{resetDebugVisualizerCamera(distance, yaw, pitch, target)},
\code{configureDebugVisualizer(flag, on)},
\code{addUserDebugLine/Text/Parameter},
\code{getKeyboardEvents()} / \code{getMouseEvents()}. The GUI has
dockable parameter panels and preview tiles. No built-in reverse.

\subsection{SOFA (C++, Qt/OpenGL GUI)}

Interactive multi-physics (deformables, biomechanics). \code{runSofa}
is an \textbf{in-process Qt app}; the default viewport is
libQGLViewer (\code{QtGLViewer}). Scenes are declarative trees
(\code{.scn} XML or Python). GUI: \textbf{Animate}, \textbf{Step},
\textbf{Reset Scene} buttons and an editable \textbf{dt} field; FPS +
elapsed time bottom-left; Save/Reset View. Mouse: left rotate, right
pan, wheel zoom, Shift+click pokes the simulation. Shortcuts:
\code{S} screenshot, \code{T} ortho/perspective, \code{V} video. No
reverse.

\subsection{Rapier (Dimforge, pure Rust)}

Modern 2-D/3-D physics engines in Rust. The core is \textbf{fully
rendering-agnostic}: \code{PhysicsPipeline::step()} advances
\code{RigidBodySet}/\code{ColliderSet}; you read poses back through
handles. The testbed (Bevy + egui) contributes the clearest control
vocabulary: \code{RunMode \{Running, Stop, Step\}}, display bitflags
(\code{TestbedStateFlags}: \code{WIREFRAME}, \code{AABBS},
\code{CONTACT\_POINTS}, \ldots), one-shot \code{TestbedActionFlags}
including \code{TAKE\_SNAPSHOT} / \code{RESTORE\_SNAPSHOT} (serialize
the whole world --- the germ of \emph{reverse}) and
\code{FRAME\_SCENE} (auto-fit camera). Its \code{GraphicsManager}
keeps hash maps from physics handles to render nodes so geometry is
sent once and only poses update.

\subsection{Open Source Physics --- OSP / EJS (Java, Davidson College)}

The classic \emph{educational} library. \code{DrawingPanel} owns the
pixel$\leftrightarrow$world transform with built-in pan/zoom;
\code{PlottingPanel}/\code{Dataset} add live graphs; \code{DataTable}
numeric readouts. \code{AbstractSimulation} runs an animation thread
calling \code{doStep()} $\sim$10$\times$/s, with a
\code{SimulationControl} GUI of \textbf{Start/Stop, Step, Reset,
Initialize} buttons, an editable parameter table and a message area;
\code{stepsPerDisplay} decouples physics rate from display rate.
Numerics live behind an \code{ODE}/\code{ODESolver} interface (Euler,
Verlet, RK4, RK45, \ldots). EJS generates Java \emph{or} JavaScript
from the same model.

\subsection{SIMple Physics (Rust, educational)}

High-school-oriented Rust simulators (SIMple Mechanics, SIMple
Gravity) built on ggez + specs ECS + imgui + Lua scripting. Editable
globals (\code{GRAVITY}, \code{PAUSED}, \code{DT}), per-object
property panels (right-click), keyboard object creation, Space pause,
\textbf{graph object properties and export to CSV} for lab reports,
Lua scene presets. Its lesson: make every quantity inspectable and
tweakable.

\subsection{PhysicsHub (web, student-built)}

Two related projects: a React/Next.js canvas app
(\code{physicshub.github.io}) and the earlier Node/p5.js ``The
Physics Hub''. The p5 version \textbf{layers canvases}:
\code{bgCanvas} (chrome + menus), \code{simCanvas} (the simulation),
\code{plotCanvas} (toggleable live plots), plus grid overlays;
sliders/buttons via p5.gui; theory text with MathJax beside every
simulation. The browser's \code{requestAnimationFrame} is the render
loop; widgets mutate parameters read by the next frame.

\subsection{The comparison chart}

\emph{(Split across two tables to stay readable; same columns.)}

\begin{center}\small
\begin{tabular}{lp{3.1cm}p{3.3cm}p{2.9cm}p{3.2cm}}
 & \textbf{PyChrono} & \textbf{PyBullet} & \textbf{SOFA} & \textbf{Rapier}\\\hline
Language / stack & C++ (+SWIG Python) & C++ engine, Python API & C++ core, Python scenes & Rust\\
Rendering & Irrlicht / VSG / OpenGL & OpenGL debug renderer & Qt + libQGLViewer & Bevy/wgpu testbed, WASM demos\\
Process model & in-process & \textbf{client--server} (shm / UDP / TCP) & in-process & in-process (lib headless)\\
Scene / camera API & \code{ChVisualSystem*}, \code{AddCamera} & \code{resetDebugVisualizerCamera}, \code{configureDebugVisualizer} & scene-graph panel, Save/Reset View & \code{GraphicsManager} maps, \code{FRAME\_SCENE}\\
Time stepping & \code{DoStepDynamics(dt)} user loop & \code{stepSimulation()}, \code{setTimeStep} & Animate / Step / dt field & \code{RunMode} \{Running, Stop, Step\}\\
Pause / step / reverse & gate the loop / --- & call-when-you-want / --- & buttons / --- & RunMode / \textbf{snapshot+restore}\\
Mouse convention & L-rotate, R-pan, wheel & L-orbit, wheel, pan & L-rotate, R-pan, wheel & orbit / pan / zoom + pick\\
Toolbar / HUD & info panel (\code{i}) & sliders, preview tiles & Animate bar, FPS+time & egui panel, stats flags\\
Async sim$\leftrightarrow$GUI & none & \textbf{command$\to$status mailbox} & none & ECS systems\\
\end{tabular}
\end{center}

\begin{center}\small
\begin{tabular}{lp{3.6cm}p{3.6cm}p{3.6cm}}
 & \textbf{OSP / EJS} & \textbf{SIMple Physics} & \textbf{PhysicsHub}\\\hline
Language / stack & Java (EJS$\to$JS too) & Rust (ggez, specs, Lua) & JS (React/Next or p5.js)\\
Rendering & Swing \code{DrawingPanel} & ggez 2-D & HTML \code{<canvas>}\\
Process model & in-process & in-process & in-browser\\
Scene / camera API & \code{DrawingPanel} pan/zoom, \code{InteractivePanel} & imgui panels, Lua globals & per-sim canvas + sliders\\
Time stepping & \code{doStep()} thread, \code{stepsPerDisplay} & \code{update()} each frame, \code{DT} global & \code{requestAnimationFrame}\\
Pause / step / reverse & Start-Stop-Step-Reset / --- & Space pause / --- & play-pause-reset / ---\\
Mouse convention & drag interactives & L-drag move, R-click edit & widget-based\\
Toolbar / HUD & buttons + param table + messages & imgui property panels & sliders, plots, theory panes\\
Async sim$\leftrightarrow$GUI & animation thread & shared ECS state & same-thread JS\\
\end{tabular}
\end{center}

\subsection{The ``scene science'' concept tree}

Which simulator taught us each concept ($\to$ = adopted in posim):

\begin{lstlisting}
scene science
|-- scene lifecycle
|   |-- create / destroy a window           <- PyBullet connect()/disconnect()  -> SCENE CREATE / CLOSE
|   |-- describe geometry once, poses often <- Rapier GraphicsManager           -> init vs frame messages
|   |-- declarative scene description       <- SOFA .scn / Bullet URDF          -> init entity list (JSON)
|   `-- refresh / redraw on demand          <- SOFA Reset Scene                 -> SCENE REFRESH / REDRAW
|-- camera control
|   |-- orbit camera (yaw,pitch,dist,target)<- PyBullet resetDebugVisualizerCamera -> SCENE ROTATE / server camera
|   |-- translate / pan                     <- Irrlicht & QGLViewer right-drag  -> SCENE TRANSLATE, arrows, shift-drag
|   |-- zoom (wheel + keys)                 <- all seven                        -> SCENE ZOOM IN/OUT/<f>, wheel, +/-
|   `-- frame-the-scene auto-fit            <- Rapier FRAME_SCENE               -> fit_distance() at create/refresh
|-- playback control
|   |-- run-mode state machine              <- Rapier RunMode                   -> Stopped/Running/Paused/Reversing
|   |-- start / stop / pause / single-step  <- OSP AnimationControl             -> SCENE START/STOP/PAUSE, steps
|   |-- reverse via snapshot history        <- Rapier TAKE/RESTORE_SNAPSHOT     -> SCENE REVERSE (ring buffer)
|   `-- settable time step                  <- PyBullet setTimeStep / SOFA dt   -> SCENE SET_TIME_STEP
|-- asynchronous messaging
|   |-- command -> status protocol          <- PyBullet client-server           -> WebSocket JSON both ways
|   |-- window -> notebook events           <- PyBullet debug/user events       -> SCENE EVENTS + {"event":...}
|   `-- transport-agnostic message structs  <- PyBullet shm/UDP/TCP             -> same JSON, any front end
`-- educational UX
    |-- toolbar of big obvious buttons      <- OSP SimulationControl            -> the toolbar
    |-- status readouts (t, E, fps)         <- SOFA FPS/time, OSP DataTable     -> the status bar
    |-- trails / plots / layered canvas     <- PhysicsHub plotCanvas            -> grid + trails + bodies layers
    |-- hide/show and visual flags          <- Bullet configureDebugVisualizer  -> SCENE HIDE / SHOW
    `-- inspect-everything philosophy       <- SIMple Physics                   -> labels, help overlay, STATUS
\end{lstlisting}

\section{Porting recommendations (and what posim actually did)}

Constraint recap: posim allows \textbf{zero external crates, zero
\code{unsafe}, zero warnings}. That immediately rules some things in
and some out.

\subsection{Recommended and adopted}

\begin{center}\small
\begin{tabular}{p{3.4cm}p{2.6cm}p{8.2cm}}
\textbf{idea} & \textbf{source} & \textbf{posim realization}\\\hline
Command$\to$status message protocol & PyBullet & Every window action is a small JSON command; the simulator answers with state. Maps 1:1 onto WebSocket text frames --- \code{std::net::TcpListener} + a hand-rolled RFC~6455 layer (SHA-1 + base64 included, $\sim$200 lines, fully unit-tested).\\
Engine/renderer decoupling & Rapier & The physics core never draws. The scene server owns a \emph{synchronized copy} of the system; the browser is just a view.\\
Geometry once, poses per frame & Rapier \code{GraphicsManager} & \code{init} message carries shapes/sizes once; 30\,Hz \code{frame} messages carry only \code{[x,y,z,qw,qx,qy,qz]} per body.\\
Run-mode state machine & Rapier \code{RunMode} & \code{Stopped / Running / Paused / Reversing} enum drives the playback thread.\\
Reverse via snapshots & Rapier snapshot / restore & A bounded ring buffer (20{,}000 frames) of cloned system states; \kw{scene reverse} replays it backward and auto-pauses at the beginning. Cheap because \code{PhysicalObjectSystem} is \code{Clone}.\\
Control-panel UX & OSP \code{SimulationControl} & Toolbar: Start/Pause/Stop/Reverse/step/dt --- the exact OSP button set plus reverse.\\
Status readouts & SOFA, OSP & Status bar: mode, t, dt, E, bodies, hidden, history, camera, fps.\\
Camera conventions & Irrlicht / QGLViewer / Bullet & Left-drag orbit, wheel zoom, right/shift-drag pan, arrows translate --- the cross-simulator standard.\\
Auto-fit camera & Rapier \code{FRAME\_SCENE} & \code{fit\_distance()} $=\max(12,\,2.5\times$ farthest entity$)$.\\
Layered drawing & PhysicsHub & Grid layer $\to$ trails layer $\to$ bodies (painter-sorted) $\to$ labels.\\
Settable dt & PyBullet \code{setTimeStep} & \kw{scene set\_time\_step} \code{<dt>} and the toolbar dt field (validated: positive, finite).\\
Debug/user events & PyBullet & Window errors and actions become notebook events (\S\ref{sec:async}).\\
\end{tabular}
\end{center}

\subsection{Considered and deliberately NOT ported}

\begin{center}\small
\begin{tabular}{p{4.6cm}p{9.6cm}}
\textbf{idea} & \textbf{why not}\\\hline
Shared-memory transport (PyBullet) & Requires raw memory mapping --- impossible without \code{unsafe} or an external crate. TCP on localhost is fast enough for 30\,Hz frames.\\
Native GUI toolkits (Qt, Irrlicht, Bevy, egui, ggez) & All are external dependencies. The browser is the one universally-available renderer; the page is embedded in the binary with \code{include\_str!}.\\
WebGL/GPU rendering & Needs no dependency, but canvas-2D with manual projection is simpler, debuggable, and fast enough for classroom scenes. A WASM/WebGL path remains open (Rapier's route).\\
Lua scripting (SIMple) & posim already has its own command language; two languages would confuse the audience.\\
Fragmented/binary WebSocket frames & Rejected explicitly by the server: text-only keeps the protocol inspectable by students.\\
\end{tabular}
\end{center}

\section{The SCENE command family (complete reference)}

All keywords are case-insensitive. Numeric arguments are
\emph{term-level} expressions: \code{-5} is negative five,
\code{2*2} works; write \code{(1+2)} to use a sum. Full grammar in
\code{grammar.pdf} \S3.

\begin{center}
\begin{longtable}{p{5.4cm}p{9.0cm}}
\textbf{command} & \textbf{what it does}\\\hline
\kw{scene create} \code{[port]} & Start the scene server (OS-assigned port if omitted), open the browser page, show all entities. Second call reports the existing URL.\\
\kw{scene close} (alias \kw{scene destroy}) & Shut the server down and disconnect every window.\\
\kw{scene translate} \code{dx dy [dz]} & Move the camera target by a world-space offset (\code{dz} defaults to 0).\\
\kw{scene rotate} \code{dyaw dpitch} & Orbit the camera: azimuth and elevation, in degrees (pitch clamps to $\pm 89^\circ$).\\
\kw{scene zoom in} / \kw{scene zoom out} & Zoom by a fixed factor (1.25$\times$).\\
\kw{scene zoom} \code{<f>} & Zoom by factor \code{f} ($>1$ zooms in; e.g.\ \code{scene zoom 0.5} zooms out).\\
\kw{scene hide} \code{[n|ALL]} & Hide object \code{objN}, or everything (default ALL).\\
\kw{scene show} \code{[n|ALL]} & Show object \code{objN}, or everything (default ALL).\\
\kw{scene refresh} & Re-copy the notebook's current system into the window (clears history).\\
\kw{scene redraw} & Re-send the full scene description to every connected window.\\
\kw{scene start} & Begin (or resume) forward time-stepped evolution.\\
\kw{scene stop} & Halt evolution and clear the recorded history.\\
\kw{scene pause} & Freeze evolution; \kw{scene start} resumes, history kept.\\
\kw{scene reverse} & Play \textbf{backward in time} through the recorded history; pauses automatically at the beginning. Errors if no history exists yet.\\
\kw{scene reset} & Re-initialize the playback: \textbf{every mutable value and the time return to their initial values} --- the state last synced at \kw{scene create}/\kw{scene refresh} --- bit-identically; history and the step counter clear and the mode returns to Stopped. \kw{scene start} (or the window's Start button) then re-runs the simulation from the beginning. The toolbar \textbf{Reset} button calls the very same primitive.\\
\kw{scene set\_time\_step} \code{dt} & Set the playback time step (must be positive and finite).\\
\kw{scene status} & Report URL, connected windows, mode, t, dt, steps, history size, entities, hidden list, camera.\\
\kw{scene events} & Drain and print the asynchronous window~$\to$~notebook event queue.\\
\end{longtable}
\end{center}

Notes for the curious:
\begin{itemize}
\item The playback loop runs on its own thread at $\sim$30 ticks/s
      (33\,ms) and advances a \textbf{synchronized copy} of your
      system --- your notebook state is never mutated by the window.
      \kw{scene refresh} re-syncs the copy; a future
      \kw{run}/\kw{step} in the notebook does not disturb a running
      scene.
\item Every forward step goes through
      \code{physical\_object::integrate} --- the same sundials
      (CVODE/ARKODE) code path as \kw{step}/\kw{run}. There is no
      second, lesser integrator hiding in the graphics.
\item \kw{reset} in the notebook keeps the window open and re-syncs
      it to the now-empty system --- including clearing the box
      wireframe and wall flags if a \kw{box} existed.
\item \kw{scene reset} is \textbf{not} the notebook's \kw{reset}:
      the notebook command empties the system itself, while
      \kw{scene reset} only rewinds the window's playback copy to the
      state it was last synced with. The reply spells it out:
      \code{scene playback reset to its initial state (t = 0, 1
      entity); Start runs the simulation again from the beginning}.
\end{itemize}

\section{The wire protocol (WebSocket, JSON text frames)}

The page connects to \code{ws://127.0.0.1:<port>/ws}. One JSON
document per frame, both directions.

\subsection{Simulator $\to$ window}

\begin{lstlisting}
{"type":"init","t":0.0,"dt":0.01,"mode":"stopped",
 "camera":{"yaw":-60.0,"pitch":55.0,"dist":12.0,"target":[0,0,0]},
 "hidden":[],
 "entities":[{"i":0,"mass":256.0,"charge":0.0,"shape":"sphere","radius":0.35},
             {"i":3,"mass":2.0,"charge":0.0,"shape":"cuboid",
              "half_extents":[0.3,0.2,0.1]}]}
\end{lstlisting}
Sent on connect, on \kw{scene refresh}/\kw{redraw}, after hide/show,
and after a playback reset (\kw{scene reset} / the Reset button).
Geometry travels \textbf{once} (the Rapier idea). Note
that \kw{box} \code{<size>} / \kw{box off} in the notebook do
\textbf{not} push a new init on their own --- box changes reach an
open window only via \kw{scene refresh}, and the \kw{box} reply
reminds you: \code{(scene window open: SCENE REFRESH shows the box)}.
The one exception is \kw{reset}, which re-syncs the window itself and
clears the box wireframe and wall flags automatically.

Each entity carries its shape-specific geometry: a \code{sphere} its
\code{radius}, a \code{cuboid} its \code{half\_extents}, a
\code{torus} its \code{ring\_radius} and \code{tube\_radius}, a
\code{disk} its \code{radius}, a \code{cylinder} its \code{radius}
and \code{half\_height}, and --- new in the dumbbell release --- a
\code{dumbbell} its two sphere radii \code{r1}/\code{r2}, its
\code{rod\_radius}, and the sphere-centre offsets
\code{z1}/\code{z2} along the local axis (the local origin is the
composite COM, so the offsets are mass-weighted and unequal). Any
entity whose object was registered with \kw{new} \code{...} \kw{as}
\code{<name>} also carries \code{"name"} --- the window prefers it
over \code{objN} for the label:

\begin{lstlisting}
{"i":6,"mass":1.0,"charge":0.0,"shape":"torus","ring_radius":1.5,"tube_radius":0.5}
{"i":9,"mass":0.6666666666666666,"charge":0.0,"shape":"disk","radius":1.0}
{"i":11,"mass":2.0,"charge":0.0,"shape":"cylinder","radius":0.5,"half_height":0.75}
{"i":0,"mass":0.0,"charge":0.0,"shape":"cuboid","half_extents":[1.0,4.0,4.0],"wall":true}
{"i":0,"mass":3.5,"charge":0.0,"shape":"dumbbell","r1":0.25,"r2":0.25,
 "rod_radius":0.1,"z1":-0.6428571428571429,"z2":0.35714285714285715,
 "name":"dumbell0"}
\end{lstlisting}

Two optional fields describe the rigid bounding box: each of the six
static wall slabs carries \code{"wall":true} on its entity (the page
then skips it in the body-drawing pass), and the message top level
gains \code{"box":4.0} --- the inner side length --- \textbf{absent}
when no box exists.

\begin{lstlisting}
{"type":"frame","t":1.234,"dt":0.002,"mode":"running","steps":617,
 "history":617,"energy":-592.49,
 "p":[0.15,0.0,0.0],"l":[0.4443,0.0,-1.506],"hidden":[],
 "contacts":[{"i":0,"j":1,"point":[0,0,0],"normal":[1,0,0],"impulse":2.0}],
 "bodies":[[x,y,z,qw,qx,qy,qz], ...]}
\end{lstlisting}
Broadcast every tick ($\sim$30\,Hz) while any window is connected.
Bodies are ordered by entity index; quaternions are \textbf{w-first}
(the posim convention everywhere). \code{contacts} carries every
collision the playback copy resolved in its latest step ---
\code{normal} is the unit action--reaction line from body \code{i}
toward body \code{j} --- and the page draws each one as a fading golden
arrow through the contact point (arrowhead on the \code{j} side,
$\sim$0.6\,s persistence; toggle with \textbf{Contacts} / \code{C}).
Verified live: a head-on elastic impact broadcasts
\code{normal [1,0,0]}, \code{point [0,0,0]}, \code{impulse 2.0} --- the
exact analytic values --- and the arrow renders at the projected
contact point.

Alongside \code{energy}, every frame now carries the other two
conserved totals, computed on the playback copy every tick: \code{p}
is the \textbf{total linear momentum} and \code{l} the \textbf{total
angular momentum about the origin} (each \code{[x,y,z]}). These are
exactly the numbers the ``conserved quantities'' readout (\S1.2)
displays.

\begin{lstlisting}
{"type":"camera","camera":{"yaw":...,"pitch":...,"dist":...,"target":[...]}}
\end{lstlisting}
Sent when the \emph{notebook} moves the camera
(\kw{scene translate}/\kw{rotate}/\kw{zoom}) so every window follows.

\subsection{Window $\to$ simulator}

\begin{lstlisting}
{"type":"cmd","action":"start"}                  // toolbar buttons:
   // start | pause | stop | reverse | reset | step | step_back | set_dt | refresh
{"type":"cmd","action":"set_dt","value":0.005}
{"type":"camera","yaw":-92,"pitch":71,"dist":7.9,"target":[0,0,0]}
   // gesture sync, throttled to 10 Hz, so SCENE STATUS matches the screen
{"type":"event","level":"error","message":"..."}  // e.g. page JS errors
{"type":"request_state"}                          // ask for a fresh init
\end{lstlisting}

Malformed messages never crash the server; they become error events
the notebook can read.

\section{Under the hood (for the reader who wants the machinery)}

\begin{lstlisting}
posim process
|-- VM thread (your commands)           SCENE ... --+
|-- listener thread (TcpListener, non-blocking poll) | Arc<Mutex<Shared>>
|-- playback thread (33 ms tick)  <------------------+
|     Running   -> history.push(system.clone()); integrate::step(dt)
|     Reversing -> system = history.pop_back()   (auto-pause at start)
|     broadcast frame JSON to every outbox (mpsc channels)
`-- per-window threads
      reader: parse WS frames -> handle_client_message (0.5 s poll for shutdown)
      writer: outbox receiver -> WS text frames (write half mutex-shared
              with the reader so pings/close never interleave a frame)
\end{lstlisting}

\begin{itemize}
\item \textbf{WebSocket layer} (\code{posim/src/scene/ws.rs}): SHA-1
      (FIPS 180-4) and base64 (RFC 4648) implemented from the
      standard, plus RFC 6455 framing (text/close/ping/pong; masked
      client frames; 126/127 length forms). Unit-tested against the
      official test vectors, including the RFC's own handshake example
      (\code{dGhlIHNhbXBsZSBub25jZQ==} $\to$
      \code{s3pPLMBiTxaQ9kYGzzhZRbK+xOo=}).
\item \textbf{HTTP layer}: \code{GET /} serves the embedded page
      (\code{include\_str!}), \code{GET /ws} upgrades. Anything else:
      404. Bound to \code{127.0.0.1} only --- the window is local by
      construction.
\item \textbf{The page} (\code{posim/src/scene/scene.html}): one
      self-contained HTML/CSS/JS file. Canvas-2D renderer with a z-up
      orbit camera, perspective projection, painter-sorted bodies,
      radial-gradient spheres, quaternion-rotated shape wireframes,
      adaptive grid, world axes, 800-point trails. No frameworks, no
      network fetches.
\item \textbf{Shape wireframes} --- every extended shape is drawn as
      a small set of loops and lines, all rotated by the body
      quaternion so spin is visible: a cuboid keeps its 12-edge
      wireframe; a \textbf{torus} is its outer and inner equators,
      two tube rings (the tube's top and bottom circles) and four
      tube cross-sections --- enough to read the hole and the spin; a
      \textbf{disk} is its rim plus two diameters; a
      \textbf{cylinder} is its two cap rims plus four side lines.
\item \textbf{The dumbbell renderer} --- a \textbf{dumbbell} is drawn
      as one rigid body: two radial-gradient-shaded spheres (radii
      \code{r1}, \code{r2}) at their COM offsets \code{z1}, \code{z2}
      rotated by the body quaternion, joined by the rod's
      \textbf{four silhouette lines} (the lines between the sphere
      centres, offset radially by \code{rod\_radius} in four
      directions). Because the offsets are mass-weighted the heavier
      end sits closer to the drawn position --- the spin and the
      asymmetry are both visible.
\item \textbf{The conserved-quantities readout} --- a permanent
      labeled overlay (\code{id="hud"}, top-left under the toolbar)
      rebuilt on every frame from the frame message's \code{energy},
      \code{p} and \code{l} fields: an \code{E} line (8 digits) and
      \code{P}/\code{L} lines showing the three components plus the
      \code{|.|} magnitude (5 digits). Entity labels, drawn in the
      body pass, prefer the init entity's \code{"name"} (the user's
      \kw{new} \code{...} \kw{as} registration) and fall back to
      \code{objN}.
\item \textbf{\code{Shared.initial} + \code{reset\_playback}} --- the
      playback state keeps an \code{initial} snapshot: the system as
      last synced from the notebook (set at \kw{scene create} and by
      every \kw{scene refresh}/\kw{reset} re-sync).
      \code{reset\_playback} restores it as a bit-identical clone ---
      every mutable value and the time return to their initial
      values --- clears the history ring and the step counter,
      returns the mode to Stopped, and flags a fresh init for every
      window. Three doors, one primitive: the toolbar \textbf{Reset}
      button, the window's \code{reset} command, and the notebook's
      \kw{scene reset} all land in the same function.
\item \textbf{The rigid bounding box} --- when the init message
      carries \code{"box"}, the page strokes the interior cube as a
      dashed \code{\#5d84a8} wireframe (12 edges, drawn just after
      the grid, beneath trails and bodies). The six static wall slabs
      that implement the box arrive tagged \code{"wall":true} and are
      skipped in the body-drawing pass: the dashed wireframe
      represents them.
\end{itemize}

\section{Asynchronous data in both directions}\label{sec:async}

The old design was strictly synchronous (one reply per request). The
scene link is genuinely asynchronous:

\textbf{Notebook $\to$ window.} Scene commands broadcast immediately;
a running playback streams frames whether or not the notebook is busy.

\textbf{Window $\to$ notebook.} The window pushes errors, data
requests and user actions at any time. They are queued (bounded at
1{,}000) and reach you two ways:

\begin{enumerate}
\item \textbf{Interactive / script mode} --- type \kw{scene events}:
\begin{lstlisting}
In [7]: scene events
Out[7]: window connected (1 total)
        window action: start
        error: canvas exploded            <- a page JS error, reported home
\end{lstlisting}
\item \textbf{\code{--machine} / Jupyter mode} --- posim emits
      \emph{unsolicited} lines
      \code{\{"event":"scene","message":"..."\}} between replies. The
      Jupyter wrapper kernel runs a background reader thread that
      routes real replies to the requester and pushes event lines
      straight to the notebook front end as \code{[scene] ...} stream
      output --- errors arrive on stderr (red), everything else on
      stdout. The \code{\{"op":"events"\}} request also drains the
      queue explicitly.
\end{enumerate}

\section{How it was verified}

\textbf{Rust tests (52 total, \code{cargo test --workspace}).} New
for the scene subsystem: SHA-1/base64/accept-key official vectors; an
HTTP test that the served page contains the toolbar, status bar and
all three gesture handlers; a full WebSocket session test (handshake
$\to$ init $\to$ camera sync $\to$ start/pause $\to$ event queue
$\to$ frame broadcast); a playback test proving
forward-then-reverse lands back \emph{exactly} on the start state;
lexer/parser/VM tests for every \kw{scene} command, including the
errors (\code{scene rotate 15} missing an argument, port $>65535$,
commands before \kw{scene create}).

\textbf{Wire-protocol test} (\code{jupyter/test\_protocol.py}, stdlib
only): 24 checks green, including the whole \kw{scene} family
end-to-end over \code{--machine} and the \code{events} op.

\textbf{Kernel test} (\code{jupyter/test\_kernel.py}, real ZMQ): 5
checks green.

\textbf{Real-browser verification} (headless Chrome 150 driven over
the DevTools protocol, dispatching genuine OS-level input events):
16/16 checks passed ---

\begin{lstlisting}
ok  scene shows all simulator entities (4)
ok  statusbar reports connection          ok  toolbar is present
ok  ArrowRight translates the view        ok  ArrowLeft translates it back
ok  ArrowUp/ArrowDown also translate
ok  left-drag rotates (yaw -92 -> -124)   ok  left-drag rotates (pitch 71 -> 87)
ok  drag does not zoom
ok  wheel up zooms in (12.00 -> 7.89)     ok  wheel down zooms out (7.89 -> 12.00)
ok  keyboard + zooms in                   ok  keyboard - zooms out
ok  toolbar Start begins evolution        ok  toolbar Pause freezes it
ok  toolbar Reverse plays backward (t 0.084 -> 0.022)
\end{lstlisting}

The gestures were then \textbf{independently re-verified} in a
second, different browser session (2026-07-22, driving genuine key,
drag and wheel events against a live 4-entity scene on port 7878):
ArrowRight moved the camera target along the view's right axis and
ArrowLeft returned it \textbf{exactly} to \code{[0, 0, 0]}; a
left-drag changed yaw $-60^\circ \to -98.8^\circ$ and pitch
$55^\circ \to 32.2^\circ$ while distance and target stayed untouched
(a pure orbit); wheel-up/down zoomed by the exact 1.15 factors
($12 \to 10.435 \to 12$) and the \code{+}/\code{-} keys by the exact
1.2 factors ($12 \to 10 \to 12$); toolbar Start advanced $t$ through
the CVODE playback path, Pause froze it, and Reverse replayed history
backward until it auto-paused at \textbf{exactly} $t = 0$ with the
initial energy $E = -194.84952$ restored to the last digit.

The four self-checking physics examples (Kepler, outer solar system,
tumbling body, charged particle) still print SUCCESS, the build is
warning-free, and \code{Cargo.lock} still lists exactly the five
local crates --- the zero-dependency proof.

\textbf{Shapes-and-box release (2026-07-24).} Two further layers of
checks accompany the \kw{torus}/\kw{disk}/\kw{cylinder} shapes and
the \kw{box} command (the workspace suite is now 94 tests: 37 lib +
15 collision + 9 conservation + 33 posim, zero warnings):

\begin{itemize}
\item \textbf{The init message carries the new geometry.} The new
      posim test
      \code{box\_shapes\_and\_wall\_flags\_reach\_the\_init\_message}
      opens a genuine WebSocket to the scene server and asserts the
      init JSON contains \code{"box":4.0}, \code{"shape":"torus"}
      with \code{"ring\_radius":1.5} and \code{"tube\_radius":0.5},
      and \code{"wall":true} on the static slab.
\item \textbf{Live browser session.} The 12-entity demo
      (\code{scripts/collisions/11\_box\_of\_shapes.posim}: six walls
      plus a torus, point, sphere, disk, cube and cylinder inside
      \kw{box} \code{4}) was played back in a real browser window:
      the window's playback copy conserved energy to
      $\sim$\code{1e-9}; a pixel scan of the canvas found the dashed
      box wireframe (5{,}904 pixels of \code{\#5d84a8}) and the
      golden contact arrows (2{,}218 pixels); toolbar Start and
      Reverse were exercised through real DOM clicks; the JS console
      stayed free of errors.
\end{itemize}

\textbf{Reset + conserved-quantities release (2026-07-24).} The
workspace suite is now \textbf{104 tests: 40 lib + 16 collision + 9
conservation + 39 posim}, zero warnings. What the new layers pin
down:

\begin{itemize}
\item \textbf{Reset is a true re-initialization.} The posim test
      \code{reset\_restores\_the\_initial\_state\_and\_start\_reruns}
      runs the playback forward ($\geq 10$ steps), calls the reset
      primitive, and asserts the mode is Stopped, the time is back to
      its initial value, the packed 13-per-object state is
      \textbf{bit-identical} to the initial pack, and the step
      counter is 0 --- then sets Running again and asserts the
      simulation re-starts. The page-serve test additionally asserts
      the permanent \code{id="bt-reset"} button is in the served
      HTML, and \code{scene\_commands\_compile} covers
      \kw{scene reset}.
\item \textbf{Reset verified live.} In a real browser session the
      demo was run to $t = 0.72$, the toolbar Reset was clicked: mode
      \code{stopped}, $t = 0$, history $= 0$, and the body's state
      bit-identically back at its start values; clicking Start then
      set the window running again from the beginning.
\item \textbf{The protocol carries the new fields.} The WebSocket
      session test (\code{websocket\_session\_end\_to\_end}) asserts
      every broadcast frame contains \code{"p":[} and \code{"l":[};
      \code{box\_shapes\_and\_wall\_flags\_reach\_the\_init\_message}
      asserts the init JSON carries the registered user name
      (\code{"name":"ringo"}).
\item \textbf{The dumbbell demo, live.} The two-dumbbell collision
      demo (\code{scripts/collisions/12\_two\_dumbbells.posim} ---
      two user-named dumbbells built by a user-defined function,
      colliding off-center with spin) was played back in a real
      browser window: the entity labels read \code{dumbell0} and
      \code{dumbell1}, and the conserved-quantities readout showed
\begin{lstlisting}
E = 7.86531061, P = [0.15000, 0.00000, 0.00000] |.| = 0.15000, L = [0.44430, 0.00000, -1.50600] |.| = 1.57017
\end{lstlisting}
      \textbf{identically before and after the impact} (L's $y$
      component showed \code{-1.31228e-13} after ---
      machine-epsilon noise). The notebook run of the same script
      reports, through 2 real CVODE collision events,
\begin{lstlisting}
t = 3 (3861 solver steps, 60 snapshots, |dE/E| = 9.463e-11, 2 collision(s) — CONTACTS lists them)
\end{lstlisting}
      with the total momentum \code{[0.15000000000000036, 0, 0]}
      \textbf{bit-identical} before and after.
\end{itemize}

\section{FAQ / troubleshooting}

\textbf{No window appeared after \kw{scene create}.}\\
The command prints the URL (e.g.\ \code{http://127.0.0.1:41234/}).
Open it in any browser yourself --- the automatic opening uses
\code{xdg-open}, which headless machines lack.

\textbf{I want a predictable port.}\\
\kw{scene create} \code{7878} $\to$ the URL is always
\code{http://127.0.0.1:7878/}. If the port is busy you get a clear
error; pick another or use \code{0} (OS-assigned).

\textbf{Stop the browser from opening (tests, servers).}\\
Set the environment variable \code{POSIM\_NO\_BROWSER=1}.

\textbf{\kw{reverse} says there is nothing to reverse.}\\
Reverse replays \emph{recorded} history. Run \kw{scene start} (or
single steps) first; \kw{stop} clears the recording, \kw{pause}
keeps it.

\textbf{The window shows old objects.}\\
The window views a synchronized \emph{copy}. After creating/deleting
objects in the notebook, type \kw{scene refresh}.

\textbf{Is any of this sent over the internet?}\\
No. The server binds \code{127.0.0.1} (loopback) only.

\textbf{Where do error messages from the window go?}\\
Into the event queue: \kw{scene events} in the notebook,
\code{[scene] ...} lines in Jupyter, \code{\{"event":...\}} lines in
\code{--machine} mode.

\end{document}
